\section{Completing kernels with two auxiliary types}\label{sec:kernels}
The family $F_k$ adds only two auxiliaries per layer. For this case we can
understand the neighbourhood obstructions by removing certain edges from the
seed. We first introduce the homology used to compare the resulting graphs,
then derive the identities and apply them to the eleven-vertex seed.

\subsection{Homology for the comparison}
A non-zero reduced homology group excludes a graph from every dismantling
class. We give the chain conventions and the elementary facts used in this
section and in Section~\ref{sec:homology}; in particular, empty
neighbourhoods must remain visible in degree $-1$.
The clique complex \(\Cl(H)\) contains the cliques of \(H\), including
the empty face. We use \defi{augmented reduced homology over \(\F\)}. For a complex \(K\), the chain group
\(\Ct_t(K)\) has the faces of cardinality \(t+1\) as a basis when
\(t\geq0\), and
\(\Ct_{-1}(K)=\F\), generated by the empty face. Lower chain groups
are zero. The boundary of a face is the sum of its codimension-one
faces; in particular the boundary of a vertex is the empty face.
Each term occurs twice in \(\partial^2\), so this is a chain complex.

We write \(\Ht_t(K)=\ker\partial_t/\operatorname{im}\partial_{t+1}\)
and call \(K\) \defi{acyclic} if every \(\Ht_t(K)\) vanishes.
In particular,
\begin{equation}\label{hom:eq:empty}
 \Ht_{-1}(\{\varnothing\})=\F,\qquad
 \Ht_t(\{\varnothing\})=0\quad(t\ne-1).
\end{equation}
A complex with a vertex has zero homology in degree \(-1\).
Its \defi{reduced Betti number} in degree \(t\) is
\(\widetilde\beta_t(K)=\dim_{\F}\Ht_t(K)\).
Acyclicity is used as a homological condition, not as a synonym for
contractibility or for membership in any \(\D_j\).
When a homological term is applied to a graph, it refers to its clique
complex; in particular, \emph{acyclic} here does not mean that the graph is a forest.

A \defi{simplicial cone} with apex \(w\) is a complex
\(w*J\) whose faces are exactly the faces \(\sigma\) of \(J\)
and the sets \(\sigma\cup\{w\}\), where \(w\notin V(J)\).
In particular, if a graph \(H\) is a cone with apex \(w\), then
\(\Cl(H)=w*\Cl(H-w)\).

\begin{lemma}\label{hom:lem:cone}
A simplicial cone is acyclic.
\end{lemma}
\begin{proof}
For an apex \(w\), let \(h\) adjoin \(w\) to a face not containing it
and send a face containing \(w\) to zero. Direct expansion gives
\(\partial h+h\partial=\mathrm{id}\), including on the empty face.
Every cycle \(z\) is therefore the boundary of \(h(z)\).
\end{proof}

We will compare complexes by expressing them as unions of acyclic pieces.

\begin{lemma}\label{hom:lem:union}
If \(K=U\cup V\) and \(U,V\) are acyclic, then for every integer \(t\),
\[
 \Ht_t(K)\cong\Ht_{t-1}(U\cap V).
\]
The intersection may be \(\{\varnothing\}\).
\end{lemma}
\begin{proof}
For \(t\geq0\), write a cycle \(z=u+v\), with chains in \(U,V\).
Then \(\partial u=\partial v=b\) is a cycle in \(U\cap V\).
Changing the decomposition adds the same chain of the intersection
to \(u\) and \(v\), so it changes \(b\) by a boundary. If \(z\) is
changed by a boundary, decompose a filling into chains of \(U,V\);
the resulting class of \(b\) is unchanged. This defines
\([z]\mapsto[b]\).

Every cycle \(b\) of the intersection has fillings \(u\) in \(U\)
and \(v\) in \(V\), since both complexes are acyclic; then \(u+v\)
proves surjectivity. If \(b=\partial a\) in the intersection, the
chains \(u+a,v+a\) are cycles in \(U,V\) and hence boundaries.
Their sum is \(z\), proving injectivity.
For \(t=-1\), the left side is zero since the acyclic complexes
\(U,V\) have vertices; the right side is in degree \(-2\).
In lower degrees both sides vanish.
\end{proof}

The folds used in the graph reductions preserve these homology groups.

\begin{lemma}\label{hom:lem:fold}
A fold \(v\to w\) induces isomorphisms
\(\Ht_t(\Cl(H-v))\cong\Ht_t(\Cl(H))\) for every integer \(t\).
\end{lemma}
\begin{proof}
Pair every face \(\sigma\) containing \(v\) but not \(w\) with
\(\tau=\sigma\cup\{w\}\). The domination condition guarantees that
\(\tau\) is a face. Remove these pairs in decreasing order of
\(|\sigma|\). When a pair is removed, \(\tau\) is the only proper
coface of \(\sigma\) still present: every larger coface was removed
in an earlier pair. The remaining faces form a subcomplex.

To see that removing one such pair preserves augmented homology,
let \(K'\) be the resulting subcomplex of the current complex \(K\)
and put \(r=\dim\sigma\). A cycle of degree \(r+1\) has coefficient
zero at \(\tau\), since its boundary has coefficient zero at
\(\sigma\). For a cycle of degree \(r\), adding its coefficient at
\(\sigma\) times \(\partial\tau\) removes \(\sigma\) without changing
its homology class. In other degrees no basis element is removed.
Thus every class has a representative in \(K'\).

Conversely, suppose a cycle \(z'\) in \(K'\) bounds a chain \(y\)
in \(K\). If \(y\) has degree \(r+1\), its coefficient at \(\tau\)
is zero, because the coefficient of \(\sigma\) in \(z'\) is zero.
If \(y\) has degree \(r\), add its coefficient at \(\sigma\) times
\(\partial\tau\); this removes \(\sigma\) and preserves its
boundary. The resulting filling lies in \(K'\). This proves
injectivity as well, including in degree \(-1\).
After all pairs are removed the complex is exactly \(\Cl(H-v)\).
\end{proof}

The recursive neighbourhood condition therefore gives a homological obstruction at every level.

\begin{lemma}\label{hom:lem:acyclic}
If \(H\in\D_j\), \(j\geq-1\), then \(\Cl(H)\) is acyclic.
\end{lemma}
\begin{proof}
Cones are covered by Lemma~\ref{hom:lem:cone}. Induct on \(j\) and,
at a fixed level, on the number of vertices. The one-vertex graph
is a cone. For a valid first deletion \(v\), the two induction
hypotheses give acyclicity of \(\Cl(H-v)\) and \(\Cl(L_H(v))\).
The exact decomposition
\[
 \Cl(H)=\Cl(H-v)\cup\bigl(v*\Cl(L_H(v))\bigr)
\]
has acyclic pieces and intersection \(\Cl(L_H(v))\).
Lemma~\ref{hom:lem:union} proves the assertion.
\end{proof}



\subsection{The extension and its kernel}
In a construction with two auxiliary types joined by an edge, the profiles
determine which seed edges contribute to the homological obstructions.
Let \(A\) be a non-empty graph, and choose \(S,T\subseteq V(A)\).
Form \(B\) by adding vertices \(p,x,y\). The centre \(p\) is adjacent
to all of \(A\) and to neither auxiliary vertex. The neighbours of \(x\) in \(A\) are exactly \(S\), those of \(y\) are exactly \(T\), and \(xy\) is an edge.
The induced graph on \(A\) remains unchanged.

\begin{lemma}\label{ker:lem:cover}
If every non-central neighbourhood of \(B\) has non-trivial
reduced homology, then \(S\cup T=V(A)\).
The same conclusion holds if \(q(L_B(v))\ne0\) for every
\(v\ne p\).
\end{lemma}
\begin{proof}
If \(a\notin S\cup T\), then its neighbours in \(B\) are \(N_A(a)\)
and \(p\). The centre is an apex of \(L_B(a)\).
This cone has zero reduced homology and \(q=0\), contradicting
either condition.
\end{proof}

Henceforth assume \(S\cup T=V(A)\), and put
\[
 S_0=S\setminus T,\qquad T_0=T\setminus S,\qquad Z=S\cap T.
\]
Define \(H\) by deleting from \(A\) all its edges between \(S_0\)
and \(T_0\), and no others. We call \(H\) the \defi{kernel of these profile data}. This name describes an operation on edges; it does
not impose a homological property on \(H\).

Set
\[
 K=\Cl(A),\quad X=\Cl(A[S]),\quad Y=\Cl(A[T]),\quad
 W=\Cl(A[Z]),\quad U=\Cl(H).
\]
Every clique of \(H\) is contained in \(S\) or in \(T\), since it
cannot meet both \(S_0\) and \(T_0\). Conversely every clique contained
in \(S\) or in \(T\) survives the edge deletions. Therefore
\begin{equation}\label{ker:eq:kernel}
 U=X\cup Y,\qquad W=X\cap Y.
\end{equation}
In particular this union is a clique complex, even if it is smaller
than \(K\).

\subsection{The two auxiliary stars}

To compare the seed, its extension and the remainder, we group the faces containing auxiliary vertices. For a vertex \(z\) outside a complex \(J\), \(z*J\) denotes its cone.
Define the subcomplex
\begin{equation}\label{ker:eq:V}
 V=(x*X)\cup\bigl(y*(Y\cup(x*W))\bigr).
\end{equation}
It consists of all faces using an auxiliary vertex together with
their subfaces. Its two displayed pieces are cones and their
intersection is \(x*W\), also a cone. This intersection contains
the vertex \(x\), even when \(W=\{\varnothing\}\).
By Lemma~\ref{hom:lem:union}, \(V\) is acyclic.
The faces of \(V\) on the original vertices are exactly \(X\cup Y\).
Consequently
\begin{equation}\label{ker:eq:decompositions}
 \Cl(B)=(p*K)\cup V,\qquad
 \Cl(B-p)=K\cup V,\qquad
 (p*K)\cap V=K\cap V=U.
\end{equation}
These are equalities of complexes.

\begin{proposition}\label{ker:prop:global}
For every integer \(t\),
\[
 \Ht_t(\Cl(B))\cong\Ht_{t-1}(U).
\]
If \(K\) is acyclic, then also
\[
 \Ht_t(\Cl(B-p))\cong\Ht_{t-1}(U).
\]
Under this extra hypothesis, \(\Cl(B-p)\) is acyclic if and only if
\(U\) is acyclic.
\end{proposition}
\begin{proof}
The two pieces in the first decomposition of
\eqref{ker:eq:decompositions} are acyclic and intersect in \(U\).
Apply Lemma~\ref{hom:lem:union}. For the second decomposition the same
argument applies when \(K\) is acyclic. The shift covers every
possible non-zero degree; hence the last equivalence follows.
\end{proof}

For a seed \(A\in\D_s\), membership of \(B-p\) in a dismantling
class therefore implies acyclicity of the kernel. The converse
statement about membership does not follow: the proposition
characterises acyclicity of the remainder, not its dismantling level.

\subsection{Original and auxiliary neighbourhoods}

The global decomposition also applies within each neighbourhood. For an original vertex, its intersection is the neighbourhood complex in the kernel.

\begin{theorem}\label{ker:thm:original-links}
For every \(a\in V(A)\) and every integer \(t\),
\[
 \Ht_t(\Cl(L_B(a)))
 \cong \Ht_{t-1}(\Cl(L_H(a))).
\]
\end{theorem}
\begin{proof}
Put \(K_a=\Cl(A[N_A(a)])\).
If \(a\in S_0\), its neighbourhood in \(B\) contains \(p,x\) but
not \(y\), and
\[
 \Cl(L_B(a))=(p*K_a)\cup(x*X_a),\qquad
 X_a=\Cl(A[N_A(a)\cap S]).
\]
Both pieces are cones and their intersection is \(X_a\).
The edge deletions defining \(H\) remove exactly the neighbours
of \(a\) in \(T_0\). Edges among its remaining neighbours in \(S\)
are unchanged. Thus \(X_a=\Cl(L_H(a))\), and
Lemma~\ref{hom:lem:union} proves the assertion. The case \(a\in T_0\)
is obtained by exchanging the two profiles and auxiliaries.

If \(a\in Z\), none of the edges incident to \(a\) is deleted.
The neighbourhood contains both auxiliaries. Restrict the
construction \eqref{ker:eq:V} to \(N_A(a)\), writing
\[
 X_a=\Cl(A[N_A(a)\cap S]),\quad
 Y_a=\Cl(A[N_A(a)\cap T]),\quad
 W_a=X_a\cap Y_a .
\]
The resulting complex \(V_a\) is acyclic by the same union-of-cones
argument. Now \(\Cl(L_B(a))=(p*K_a)\cup V_a\), with intersection
\(X_a\cup Y_a\). The proof of \eqref{ker:eq:kernel}, restricted to
\(N_A(a)\), identifies this intersection with \(\Cl(L_H(a))\).
Lemma~\ref{hom:lem:union} again applies.
Augmented chains cover all possible empty neighbourhoods.
\end{proof}

The auxiliary neighbourhoods have the exact descriptions
\begin{equation}\label{ker:eq:auxiliary-complexes}
 \Cl(L_B(x))=X\cup(y*W),\qquad
 \Cl(L_B(y))=Y\cup(x*W).
\end{equation}
They do not contain \(p\).

\begin{lemma}\label{ker:lem:intersection-map}
If \(U=X\cup Y\) is acyclic, the inclusions induce isomorphisms
\[
 \Ht_t(W)\ \longrightarrow\ \Ht_t(X)\oplus\Ht_t(Y)
\]
for all integers \(t\).
\end{lemma}
\begin{proof}
For surjectivity, let \(a,b\) be cycles in \(X,Y\).
Since \(a+b\) is a cycle in \(U\), it bounds a chain
\(w=w_X+w_Y\) there. Then
\(z=a+\partial w_X=b+\partial w_Y\) is a cycle in \(W\)
representing the prescribed pair of homology classes.
For injectivity, suppose a cycle \(z\) in \(W\) bounds chains
\(u\) in \(X\) and \(v\) in \(Y\). The cycle \(u+v\) bounds
\(w=w_X+w_Y\) in \(U\). The chain
\(u+\partial w_X=v+\partial w_Y\) lies in \(W\) and has
boundary \(z\). These arguments apply in degree \(-1\) as well,
using the augmentation, and lower degrees are zero.
\end{proof}

When the kernel is acyclic, the auxiliary neighbourhoods are controlled by the opposite profile subgraphs.

\begin{theorem}\label{ker:thm:auxiliary-links}
If \(U\) is acyclic, then for every integer \(t\),
\[
 \Ht_t(\Cl(L_B(x)))\cong\Ht_{t-1}(Y),\qquad
 \Ht_t(\Cl(L_B(y)))\cong\Ht_{t-1}(X).
\]
\end{theorem}
\begin{proof}
By Lemma~\ref{ker:lem:intersection-map}, the map
\(\Ht_r(W)\to\Ht_r(X)\) is surjective and its kernel is
isomorphic to \(\Ht_r(Y)\): under the lemma's isomorphism
it is the projection onto the first summand.

For \(t\geq0\), write a cycle of \(X\cup(y*W)\) as \(u+y*w\).
The equations for being a cycle are
\(\partial w=0\) and \(\partial u=w\).
Assign to it the class of \(w\), which lies in the kernel of
\(\Ht_{t-1}(W)\to\Ht_{t-1}(X)\).
The assignment respects boundaries: the boundary of \(a+y*b\)
changes \(w\) by \(\partial b\).
Every class in that kernel has a filling \(u\) in \(X\), so
the assignment is surjective.

If \(w=\partial d\) in \(W\), adding
\(\partial(y*d)=d+y*\partial d\) leaves the cycle \(u+d\) in \(X\).
Surjectivity of \(\Ht_t(W)\to\Ht_t(X)\) lets us replace this
cycle, modulo a boundary in \(X\), by a cycle in \(W\).
That cycle bounds in the cone \(y*W\).
The original cycle therefore bounds in \(X\cup(y*W)\),
proving injectivity. Its homology in degree \(-1\) is zero
because it contains \(y\); the claimed right side then has
degree \(-2\) and is also zero. Lower degrees vanish.
Use \eqref{ker:eq:auxiliary-complexes} and exchange \(X,Y\)
for the second identity.
\end{proof}

\subsection{Equivalent finite conditions and their Euler form}

Here the layer construction of Section~\ref{sec:construction} has
$d=2$, $R=K_2$ and profiles $S,T$. We may therefore write the auxiliaries
as $x_i,y_i$, and $|V(G_m)|=|V(A)|+3m$. The identities just proved separate
the obstructions on original vertices from those on the two auxiliary types.

\begin{theorem}\label{ker:thm:equivalence}
Assume \(S\cup T=V(A)\), \(s\geq0\),
\(A\in\D_s\setminus\D_{s-1}\), and \(B-p\in\D_{s+1}\).
Then the following are equivalent:
\begin{enumerate}
\item Every non-central neighbourhood \(L_B(v)\), \(v\ne p\),
has non-trivial reduced homology.
\item The clique complex of the kernel \(H\) is acyclic, every \(L_H(a)\),
\(a\in V(A)\), has non-trivial reduced homology, and both
\(A[S]\) and \(A[T]\) have non-trivial reduced homology.
\end{enumerate}
Either condition implies
\[
 G_m\in\D_{s+m}\setminus\D_{s+m-1}\quad(m\geq0).
\]
\end{theorem}
\begin{proof}
The seed and remainder hypotheses imply acyclicity of \(K\) and
\(\Cl(B-p)\). Proposition~\ref{ker:prop:global} then implies
acyclicity of \(U=\Cl(H)\).
Theorem~\ref{ker:thm:original-links} gives equivalence of the
conditions on original vertices. Theorem~\ref{ker:thm:auxiliary-links}
gives equivalence of those on the auxiliaries.
Under condition (1), the data satisfy precisely the homological
criterion in Theorem~\ref{hom:thm:criterion}, proved independently in
Section~\ref{sec:homology}. It applies to
the displayed layered construction and proves strictness at
every subsequent level.
\end{proof}

The remainder hypothesis is present on both sides of this
equivalence. Acyclicity of the kernel is redundant under that
hypothesis, but states explicitly the property needed when
starting with a kernel. The equivalence does not dispense with the
remainder hypothesis.
Lemma~\ref{ker:lem:cover} gives the covering condition on the profiles.

Extend \(q\) from graphs to finite complexes by putting
\(q(J)=\sum_{\sigma\in J}(-1)^{|\sigma|}\), including the empty
face. This agrees with the existing notation \(q(H)=q(\Cl(H))\).
Inclusion--exclusion and the zero value on cones give identities
without homological assumptions.

\begin{proposition}\label{ker:prop:euler}
For the profile data above,
\begin{align}
 q(B)&=-q(H),&
 q(B-p)&=q(A)-q(H),\label{ker:eq:global-q}\\
 q(L_B(a))&=-q(L_H(a)) &&(a\in V(A)),\label{ker:eq:fixed-q}\\
 q(L_B(x))&=q(A[S])-q(A[Z]),&
 q(L_B(y))&=q(A[T])-q(A[Z]).\label{ker:eq:aux-q}
\end{align}
If \(q(H)=0\), the last line becomes
\[
 q(L_B(x))=-q(A[T]),\qquad
 q(L_B(y))=-q(A[S]).
\]
\end{proposition}
\begin{proof}
In \eqref{ker:eq:V}, both pieces and their intersection are cones,
so \(q(V)=0\). Apply inclusion--exclusion to
\eqref{ker:eq:decompositions} to obtain \eqref{ker:eq:global-q}.
The decompositions used to prove
Theorem~\ref{ker:thm:original-links} likewise have pieces of
zero \(q\) and intersection \(\Cl(L_H(a))\), yielding
\eqref{ker:eq:fixed-q}.
Apply inclusion--exclusion to \eqref{ker:eq:auxiliary-complexes}
to obtain \eqref{ker:eq:aux-q}.
Finally \eqref{ker:eq:kernel} gives
\(q(H)=q(A[S])+q(A[T])-q(A[Z])\).
\end{proof}

In particular, under \(A\in\D_s\setminus\D_{s-1}\) and
\(B-p\in\D_s\), both have \(q=0\), so \(q(H)=0\).
The non-zero Euler conditions on all non-central neighbourhoods
of \(B\) are then equivalent to
\[
 q(L_H(a))\ne0\quad\hbox{for all }a\in V(A),\qquad
 q(A[S])\ne0,\quad q(A[T])\ne0.
\]
This is the kernel formulation of Theorem~\ref{thm:template}. The identities remain valid if the remainder
is only assumed to lie in \(\D_{s+1}\), as in its separate
relaxation. The equality \(q=0\) is not used as a
test for acyclicity or dismantlability.

One further necessary condition can be read from
Lemma~\ref{ker:lem:intersection-map}. If \(S,T,Z\) are non-empty
and \(H\) is acyclic, the numbers of connected components satisfy
\[
 c(A[Z])=c(A[S])+c(A[T])-1.
\]
Indeed, take dimensions in degree zero and use
\(\dim\Ht_0=c-1\).
This numerical equality alone does not imply acyclicity; even
agreement of dimensions in all degrees would not establish that
the inclusion maps in the lemma are isomorphisms.

\subsection{Completing one edge and reducing the remainder}

The case of a single added edge lets us give an explicit reduction of the remainder. Suppose \(H\) is a graph, \(ab\) is a non-edge, and
\[
 A=H+ab,\qquad S=V(H)\setminus\{b\},\qquad
 T=V(H)\setminus\{a\}.
\]
Here \(S_0=\{a\}\) and \(T_0=\{b\}\), so deleting the edges between
the exclusive profile parts recovers exactly \(H\).

\begin{lemma}\label{ker:lem:delete-vertex}
If \(\Cl(H)\) is acyclic, inclusion induces
\[
 \Ht_t(\Cl(L_H(v)))\cong\Ht_t(\Cl(H-v))
\]
for every vertex \(v\) and every integer \(t\).
\end{lemma}
\begin{proof}
Write \(P=\Cl(H-v)\), \(Q=v*\Cl(L_H(v))\) and
\(I=P\cap Q=\Cl(L_H(v))\). The union is acyclic by hypothesis,
and \(Q\) is a cone. The same chain argument as in
Lemma~\ref{ker:lem:intersection-map}, now with the acyclic union
\(P\cup Q=\Cl(H)\), shows that
\(\Ht_t(I)\to\Ht_t(P)\oplus\Ht_t(Q)\) is an isomorphism.
The second summand is zero.
\end{proof}

Thus, if \(H\) is acyclic and all its neighbourhoods have
non-trivial reduced homology, then both \(A[S]=H-b\) and
\(A[T]=H-a\) have non-trivial reduced homology.
The remaining obligations for
Theorem~\ref{ker:thm:equivalence} are the strict level of the completed
seed and the required level of \(B-p\).

There is a useful sufficient condition for the latter.
For disjoint graphs \(C,J\), write \(C\star J\) for their
\defi{complete join}.

\begin{lemma}\label{ker:lem:join}
If \(C\in\D_j\), \(j\geq-1\), then
\(C\star J\in\D_j\) for every graph \(J\).
\end{lemma}
\begin{proof}
For \(j=-1\), an apex of \(C\) remains universal.
Induct on \(j\). Follow a level-\(j\) reduction of \(C\).
The neighbourhood of a deleted vertex in the join is the join
of its current neighbourhood in \(C\) with \(J\), and belongs
to \(\D_{j-1}\) by induction. When only one vertex of \(C\)
remains, the join is a cone and reduces at level zero, hence at
level \(j\). This proves the assertion for \(j\geq0\).
\end{proof}

\begin{proposition}\label{ker:prop:remainder-folds}
For the one-edge completion, put
\[
 C=H[N_H(a)\cap N_H(b)].
\]
The remainder \(B-p\) folds onto \(C\star C_4\), where the cycle
has vertices \(a,b,y,x\) in this order.
Consequently, if \(C\in\D_j\) for some \(j\geq0\), then
\(B-p\in\D_j\).
\end{proposition}
\begin{proof}
Every \(z\in V(H)\setminus\{a,b\}\) that is not adjacent to \(a\) in \(H\)
is dominated by \(y\) in \(B-p\). Indeed \(y\) is adjacent to
every vertex except \(a\), and \(a\) is not in the closed
neighbourhood of \(z\). Likewise, a vertex not adjacent to \(b\)
is dominated by \(x\).
Delete all vertices outside \(C\cup\{a,b,x,y\}\), retaining
\(a,b,x,y\) and the common neighbours. Each chosen domination
inclusion remains true after restricting to the vertices still
present.

In the remaining graph, the edges on \(\{a,b,x,y\}\) are exactly
\(ab,by,yx,xa\), and every vertex of \(C\) is adjacent to these
four vertices. Its induced graph is unchanged from \(H\).
The remainder is therefore \(C\star C_4\); when \(C\) is empty,
it is just the four-cycle.
If \(C\in\D_j\), Lemma~\ref{ker:lem:join} gives
\(C\star C_4\in\D_j\). Prepending the displayed folds is valid
at level \(j\) by nesting of the classes.
\end{proof}

If \(C\) is a cone, the proposition applies with \(j=0\).
The argument gives a sufficient certificate of the remainder.
It does not assert that a failure to dismantle \(C\) at a chosen
level excludes membership of \(B-p\) at that level.
In particular, an unsuccessful sufficient test must not replace
the original remainder condition in the equivalence theorem.

\subsection{The completion underlying the eleven-vertex seed}

Define the completed seed \(A\) on \(0,\ldots,10\) by
Table~\ref{ker:tab:seed}. It has thirty-three edges.
The pair \(\{1,2\}\) is an edge; define
\(V(H)=V(A)\) and \(E(H)=E(A)\setminus\{\{1,2\}\}\).
Then \(H\) has thirty-two edges and eleven vertices.
For the profiles \(S=V(A)\setminus\{2\}\) and
\(T=V(A)\setminus\{1\}\), the exclusive parts are
\(\{1\}\) and \(\{2\}\), respectively. Hence this \(H\) is
exactly the kernel of the profile data, and adding the single edge
\(\{1,2\}\) to \(H\) gives \(A\).

\begin{table}[htbp]
\centering
\caption{The complete adjacency data of the completed seed \(A\).}
\label{ker:tab:seed}\label{tab:seed-adjacency}
\begin{tabular}{cl}
\toprule
\(v\)&\(N_A(v)\)\\\midrule
$0$ & $\{3,4,6,7,8\}$\\
$1$ & $\{2,4,5,6,7,9,10\}$\\
$2$ & $\{1,5,6,7,8,9,10\}$\\
$3$ & $\{0,6,7,8,9,10\}$\\
$4$ & $\{0,1,7,8,9,10\}$\\
$5$ & $\{1,2,6,8,9\}$\\
$6$ & $\{0,1,2,3,5,10\}$\\
$7$ & $\{0,1,2,3,4,9\}$\\
$8$ & $\{0,2,3,4,5,10\}$\\
$9$ & $\{1,2,3,4,5,7\}$\\
$10$ & $\{1,2,3,4,6,8\}$\\
\bottomrule
\end{tabular}

\end{table}

This is the completed seed $F_1$ used in the article. The reductions of
Appendix~\ref{sec:article} establish its strict level and dismantle $B-p$.
The homology calculation gives
\[
 \Ht_*(\Cl(H))=0,\qquad
 \dim_{\F}\Ht_1(\Cl(L_H(v)))=
 \begin{cases}2,&v=10,\\1,&v\ne10,\end{cases}
\]
with all other reduced homology groups of these neighbourhoods zero.
Moreover, $N_H(1)\cap N_H(2)=\{5,6,7,9,10\}$ induces the path
$10,6,5,9,7$. Proposition~\ref{ker:prop:remainder-folds} therefore gives
$B-p\in\D_0$. Explicitly, the folds $0,3,8\to y$ followed by $4\to x$
leave the join of this path with the four-cycle on $1,2,y,x$.
The displayed neighbourhood groups and Lemma~\ref{ker:lem:delete-vertex}
give condition~(2) of Theorem~\ref{ker:thm:equivalence}.
With the strict level of $A$ proved in Appendix~\ref{sec:article}, that
theorem provides a homological proof of the strict levels of $F_k$.
\runcheck{kernels}
The report reconstructs $H$ by removing $\{1,2\}$, restores the seed by adding
that edge, computes the displayed homology groups over $\F$, and checks the
common-neighbour graph. It also compares the homology of each neighbourhood
in $H$ with the degree-shifted group in $B$.
